\documentclass[11pt]{article}
\input{../../shared/project_style.tex}
\rhead{Basics 6 Textbook Draft}

\begin{document}
\DocTitle{Basics 6: Linear Algebra and Eigenvalue Problems}{I - Mathematical and computational foundations}{Vector spaces, inner products, generalized eigenproblems, conditioning, and physical modes}

\begin{overviewbox}
\textbf{Textbook draft, version 1.0.} Discretized plasma equations become algebraic systems. Equilibria require solves, waves require eigenpairs, and multiphysics coupling requires stable linearization and conditioning. This chapter is designed for a reader with no prior plasma-physics training. It distinguishes exact identities from physical approximations and connects the central equations to later simulation modules.
\end{overviewbox}

\ProjectTOC

\section{Learning objectives}
After completing this note, the reader should be able to:
\begin{itemize}[leftmargin=1.6em]
\item Interpret arrays and functions as vectors in a vector space.
\item Use bases, inner products, norms, and orthogonality.
\item Explain rank, null spaces, solvability, and condition number.
\item Derive ordinary and generalized eigenvalue problems.
\item Use symmetry, positivity, and Rayleigh quotients.
\item Assess residuals, normalization, convergence, and mode tracking.
\end{itemize}

\section{Prerequisites and reading strategy}
\begin{itemize}[leftmargin=1.6em]
\item Basics 1 and elementary algebra; integration for continuous inner products.
\end{itemize}
On a first reading, emphasize physical meaning and dimensional checks. On a second reading, reproduce the derivations. Attempt the computational laboratory only after the limiting cases are understood.

\section{Vector spaces and bases}

A vector space is closed under addition and scalar multiplication. Functions form vector spaces as naturally as finite arrays. In a basis $\{\mathbf e_i\}$,
\begin{equation}
\mathbf x=\sum_i x_i\mathbf e_i.
\end{equation}
Discretization represents a continuous field by a finite coordinate vector. The basis may be grid values, finite-element functions, Fourier modes, or learned features.


\section{Inner products and orthogonality}

The Euclidean inner product is $\mathbf x^T\mathbf y$. A weighted inner product is
\begin{equation}
\langle\mathbf x,\mathbf y\rangle_M=\mathbf x^TM\mathbf y,
\end{equation}
with symmetric positive-definite $M$. A continuous analogue is
\begin{equation}
\langle\xi,\eta\rangle=\int w(s)\xi(s)\eta(s)\,ds.
\end{equation}
Orthogonality depends on the chosen metric. Generalized eigenmodes are often orthogonal in a mass-matrix metric rather than the ordinary dot product.


\section{Linear systems, null spaces, and conditioning}

A square system $A\mathbf x=\mathbf b$ has a unique solution if $A$ is nonsingular. The null space contains solutions of $A\mathbf x=0$. Gauge freedoms, pure-Neumann boundary conditions, and unconstrained normalizations often create null spaces.

The condition number
\begin{equation}
\kappa(A)=\norm{A}\norm{A^{-1}}
\end{equation}
measures worst-case amplification of input error. Nondimensionalization, scaling, and preconditioning improve conditioning without changing the intended physical solution.


\section{Eigenvalues and dynamical modes}

An eigenpair satisfies
\begin{equation}
A\mathbf x=\lambda\mathbf x,\qquad\mathbf x\ne0.
\end{equation}
For $\dot{\mathbf q}=A\mathbf q$, a modal solution $\mathbf q=\mathbf xe^{\lambda t}$ has growth and oscillation set by $\lambda$.

A second-order system
\begin{equation}
M\ddot{\mathbf q}+K\mathbf q=0
\end{equation}
with $\mathbf q=\mathbf xe^{-i\omega t}$ gives
\begin{equation}
K\mathbf x=\omega^2M\mathbf x.
\end{equation}
$K$ represents restoring forces and $M$ inertia. Module 4 uses this structure.


\section{Symmetry, positivity, and Rayleigh quotients}

If $K$ and $M$ are symmetric and $M$ is positive definite, generalized eigenvalues are real under appropriate boundary conditions. The Rayleigh quotient
\begin{equation}
\mathcal R(\mathbf x)=\frac{\mathbf x^TK\mathbf x}{\mathbf x^TM\mathbf x}
\end{equation}
equals $\omega^2$ at an eigenvector. Positive $K$ implies positive squared frequencies. An indefinite force operator may produce negative $\omega^2$, signaling ideal instability.


\section{Residuals, normalization, sparse solvers, and mode identity}

A computed eigenpair should satisfy a normalized residual such as
\begin{equation}
r=\frac{\norm{K\mathbf x-\lambda M\mathbf x}}
{\norm{K\mathbf x}+|\lambda|\norm{M\mathbf x}}.
\end{equation}
A small residual proves algebraic consistency, not continuum convergence. Eigenvectors have arbitrary scale and sign, so impose
\begin{equation}
\mathbf x^TM\mathbf x=1
\end{equation}
and a reproducible sign rule.

Sparse Krylov and shift-invert methods find selected modes without diagonalizing the whole matrix. During parameter scans, frequency sorting can swap branch identity near avoided crossings. Use metric overlaps, harmonic weights, and localization.


\section{Computational laboratory}

Assemble mass and stiffness matrices for a fixed-end mass-spring chain. Solve $Kx=\omega^2Mx$, normalize in the $M$ inner product, and verify orthogonality and residuals. Refine the chain and compare with the analytic string limit. Perturb one spring and track modes by overlap rather than eigenvalue rank.


\section{Summary}
\begin{itemize}[leftmargin=1.6em]
\item Discretization maps fields into finite-dimensional vector spaces.
\item Inner products define normalization and orthogonality.
\item Generalized eigenproblems separate inertia and restoring forces.
\item Symmetry and positive definiteness strongly constrain the spectrum.
\item Residuals, refinement, and mode tracking are all required.
\end{itemize}

\SymbolsAcronymsSection
\begin{symtable}
$A,K,M$ & operators & Generic, stiffness, and mass matrices. \\
$\lambda,\omega^2$ & eigenvalue & Modal growth parameter or squared frequency. \\
$\mathbf x$ & eigenvector & Discrete mode shape. \\
$\kappa(A)$ & condition number & Sensitivity of a linear solve. \\
$\mathcal R$ & Rayleigh quotient & Energy ratio estimating an eigenvalue. \\
\end{symtable}

\section{Exercises}
\begin{enumerate}[leftmargin=1.8em]
\item Prove orthogonality of distinct eigenvectors of a real symmetric matrix.
\item Derive the generalized Rayleigh quotient.
\item Give an example where a small matrix residual does not imply a good physical solution.
\item Explain why eigenvector scale and sign are arbitrary.
\item Derive the eigenvalues of a coupled $2\times2$ matrix.
\end{enumerate}

\section{References}
\begin{thebibliography}{99}
\bibitem{ref1} G. B. Arfken, H. J. Weber, and F. E. Harris, \emph{Mathematical Methods for Physicists}, 7th ed., Academic Press (2013).
\bibitem{ref2} G. Strang, \emph{Linear Algebra and Its Applications}, 4th ed., Brooks/Cole (2006).
\bibitem{ref3} W. A. Strauss, \emph{Partial Differential Equations: An Introduction}, 2nd ed., Wiley (2007).
\bibitem{ref4} J. P. Freidberg, \emph{Ideal MHD}, Cambridge University Press (2014).
\bibitem{ref5} J. P. Goedbloed, R. Keppens, and S. Poedts, \emph{Advanced Magnetohydrodynamics}, Cambridge University Press (2010).
\bibitem{ref6} H. Goedbloed and S. Poedts, \emph{Principles of Magnetohydrodynamics}, Cambridge University Press (2004).
\end{thebibliography}

\end{document}
